\documentclass[trsc,nonblindrev]{informs3noheader}
\newtheorem{theorem}{Theorem}
\newtheorem{lemma}{Lemma}
\usepackage[utf8]{inputenc}
\usepackage[ruled,vlined,linesnumbered]{algorithm2e}
\usepackage{verbatim}
\newcommand{\manxi}[1]{{\color{blue}[Manxi: #1]}}
\newcommand{\response}[1]{{\color{blue}[#1]}}
\usepackage[colorlinks=true,allcolors=blue]{hyperref}

\usepackage[
  backend=biber,
  style=authoryear,
  natbib=true
]{biblatex}

\addbibresource{refs.bib}


\OneAndAHalfSpacedXI

\begin{document}

\TITLE{Lime operations}

\ARTICLEAUTHORS{
\AUTHOR{}
\AFF{
\EMAIL{}}
}
\maketitle


\section{Model Description}
\begin{itemize}
    \item Region (Nodes) $i \in N$ 
    \item Finite time horizon $t=1, 2, \dots, T$
    \item Available bikes $A_{i}^t$ and unavailable bikes $U_{i}^t$
    \item Time cost for traveling from node $i$ to $j$ via scooter $t_{ij}$
    \item Worker time cost of traveling from $i$ to $j$ is $d_{ij}$
    \item Task operation cost from $i$ to $j$ is $c_{ij}$
    \item Potential demand: $D_{ij}^t$ (e.g.\ Poisson with rate $\lambda_{ij}^t$). The total demand of trips start from station $i$ is
    \[
    D_i^t={\sum_{j \in N}D_{ij}^t}
    \]
    \item Served demand: The number of demand starts from station $i$ is
    \[
    Y_i^t = \min\{A_{i}^t,\, D_i^t\}
    \]
    The destination specified demand for each destination $j$ is given by
    \[
    Y_{ij}^t = Y_i^t \times {\frac{D_{ij}^t}{D_i^t}}
    \]
   
    \item Lost demand: 
    Since the penalty for unsatisfied demand is not o-d based, we do not need to track the demand loss on each o-d. The number of unsatisfied demand from station $i$ is
    \[
      L_{i}^t = D_i^t - Y_i^t
    \]
    \item $R_{ij}$ is the revenue platform gets from customer trip $ij$, $C_p$ is the penalty for unsatisfied demand.
    
    \item A bike becomes unavailable after trip $ij$ with probability $\phi_{ij}$
    \item Bikes arriving at $j$ by users at time $t$ that remain available 
    \[F_j^t = \sum_{i \in N} Y_{ij}^{t- t_{ij}}(1-\phi_{ij})\]Bikes that arrive at $j$ at time $t$ that are unavailable 
    \[\bar{F}_{j}^t = \sum_{i \in N} Y_{ij}^{t- t_{ij}}\phi_{ij}\]
    
    \item Task posted at time $t$: 
    \begin{itemize}
    \item Number of bikes to swap at station $i$: $\hat{m}_{ii}^t$, should have
    \[
    \hat{m}_{jj}^t\leq U_j^t + \bar{F}_j^t
    \]
    \item Number of bikes to move from $i$ to $j$: $\hat{m}_{ij}^t$, should have
    \[
    \hat{m}_{ij}^t\le A_j^t - Y_j^t
    \]
     \end{itemize}
     
     \item Number of swap jobs that are open at time $t$ at node $i$: $M_{ii}^t$. Number of moving jobs that are open at time $t$ from $i$ to $j$: $M_{ij}^t$.

     \item Price that platform pay for swap task: $p_{ii}$; price that platform pay for move task from $i$ to $j$: $p_{ij}$, 
     
     \item Number of workers available in station $i$ at epoch $t$: $W_i^t$; These workers add up to $W$
     \item Number of swap jobs being chosen at station $i$ at time $t$:  $\tilde{s}_i^t$; Number of moving jobs being chosen from $i$ to $j$: $\tilde{m}_{ij}^t$

     \item The workers' flow due to task allocation is $x_{ijk}^t$, which denotes the number of workers who starts from $i$, goes to $j$ for working and ends up in $k$;

    \item Task payments and worker net profit:
    \begin{itemize}
        \item A worker $w$ located at $\ell(w)$ who undertakes a task $m$, which starts from $i$ and ends at $j$, first travels from $\ell(w)$ to $i$ and then executes the task from $i$ to $j$.
        The corresponding net profit for worker $w$ is
        \[
            \mathcal{P}_w^t(i,j) = p_{ij} - d_{\ell(w),\,i} - c_{ij},
        \]
        
        \item We assume that workers have an outside option of remaining idle with profit $0$, and they regard any task trip $(i,j)$ with $\mathcal{P}_w^t(i,j)  < 0$ as unacceptable.
    \end{itemize}

    \item Task Allocation Problem As Stable Matching

    \begin{itemize}
    \item We model a stable matching problem between workers from worker set $\mathcal{W}^t$ and task trips from task set $\mathcal{M}^t$
    
    \item We think of workers as greedy: when they are available, they choose the job with highest profit that they can get. A job is obtained by the nearest available worker.
    \end{itemize}
    
\begin{comment}
    The algorithm guarantees termination within $\mathcal{O}$($\left |\mathcal{W}^t\right|\times \left |\mathcal{M}^t\right|$) proposals; 
    
    The algorithm guarantees Pareto Optimality along workers within all stable matching\parencite{gale1962college}

    % {\color{red}The algorithm shown above is equivalent to solving the following Integer Programming \parencite{agoston2016integer}:
\end{comment}
     \item Platform chooses:

\begin{itemize}
    \item Contract prices $p_{ij}$, where $i,j\in \mathcal{N}$
    \item Task posting policy $\hat{m}_{ij}^t$ 
    
    % (mapping task state $\mathcal{S}$, bike state $A_i^t$, $U_i^t$) 
    %\[
    %\pi\big[( \,(s_i)_{i\in N},(m_{ij})_{i,j\in N} \,)\mid\mathcal{S}, (A_i)_{i\in N}, (U_i)_{i,j\in N}, \,(W_i)_{i\in N}\big]
    %\]
    \begin{itemize}
        \item Optimization aspect: for each epoch $t$, choose optimal number of task post: $\hat{m}_{ij}^t$, where $i,j\in N$, to maximize expected revenue
    \end{itemize}

\end{itemize}

\newpage

\item Problem Formulation 
\begin{subequations}
\begin{align}
    \textbf{objective:}\nonumber\\
    \mathbf{\max} \quad
    &\mathbb{E}\left[
    \sum_{t=1}^{T}
    \left(
    \sum_{i\in \mathcal{N}} \big( \sum_{j\in \mathcal{N}}R_{ij}\, Y_{ij}^t - C_p\, L_{i}^t \big)
    - {\color{red}\sum_{j,k\in \mathcal{N}} p_{jk}\, \tilde{m}_{jk}^t}
    \right)
    \right] \label{obj}\\[4pt]
    \textbf{subject to:}\nonumber\\
    & {\color{red} Y_i^t \leq A_{i}^t} & \forall i\in \mathcal{N}, t\in [T] \label{min(A,D)start}\\[4pt]
    & {\color{red} Y_i^t \leq D_{i}^t} & \forall i\in \mathcal{N}; t\in [T]\\[4pt]
    & {\color{red}Y_i^t \geq  \alpha_i A_{i}^t + (1-\alpha_i)  D_{i}^t} & \forall i\in \mathcal{N}; t\in [T] \label{min(A,D)}\\[4pt]
    & {\color{red} \alpha_i \leq 1} & \forall i\in \mathcal{N} \\
    & Y_{ij}^t = Y_i^t \times {\frac{D_{ij}^t}{D_i^t}}   &\forall i,j \in \mathcal{N}; t \in [T] \label{OrderFlow}\\[4pt]
    & L_{i}^t = D_i^t - Y_i^t &\forall i \in \mathcal{N}; t \in [T] \label{OrderLoss}\\[4pt]
    & F_j^t = \sum_{i \in \mathcal{N}} Y_{ij}^{t- c_{ij}}(1-\phi_{ij}) &\forall  j \in \mathcal{N};t \in [T] \label{ReturningAvailable}\\[4pt] 
    %%%%% above is constraint 51 %%%%%
    & \bar{F_{j}^t} = \sum_{i \in \mathcal{N}} Y_{ij}^{t- c_{ij}} \phi_{ij} &\forall j \in \mathcal{N};t \in [T] \label{ReturningUnavailable}\\[4pt]
    & \hat{m}_{jj}^t\leq U_j^t + \bar{F}_j^t &\forall  j \in \mathcal{N};t \in [T] \label{SwapConstraint}\\[4pt]
    & \hat{m}_{ij}^t\le A_i^t - Y_i^t &\forall  i,j \in \mathcal{N},i\neq j; t \in [T] \label{MoveConstraint}\\[4pt]
    &A_j^{t+1}
            = A_j^{t}
            - Y_j^{t} % Realized services
            + F_j^{t} % Returning bike that is not battery dead (Services completed)
            - \sum_{k\in \mathcal{N}\setminus \{j\} } \hat{m}_{jk}^{t} % Tasks posted except for swap task, since swap task is tracked by F_j
            + \sum_{i,k\in \mathcal{N}} x_{ikj}^{t-d_{ik}-c_{kj}} % all tasks that is to be finished at time t
            &\quad \forall j \in \mathcal{N};\;t\in[T]\label{FlowTransition}\\[4pt]
    &U_j^{t+1} % Only counts battery failure
        = U_j^{t}
        + \bar{F}_j^{t} % Returning dead bikes
        - \sum_{i\in \mathcal{N}}x_{ijj}^{t-d_{ij}-c_{jj}} % 
            &  \quad \forall j \in \mathcal{N}; \;t\in[T] \\[4pt]
    &\sum_{j,k\in \mathcal{N}} x_{ijk}^t \le 
        W_i^t, &\quad \forall i\in \mathcal{N};\; t\in [T]\\[4pt]
    &W_k^{t+1}
        = W_k^t
        - \sum_{i,j\in \mathcal{N}} x_{kij}^{t}
        + \sum_{i,j\in \mathcal{N}} x_{ijk}^{t-d_{ij}-c_{jk}}
       &\quad \forall k\in \mathcal{N};\; t\in [T]\\[4pt]
    &M_{ij}^{t+1}
        = M_{ij}^t - \tilde{m}_{ij}^{t} + \hat{m}_{ij}^{t+1},
        & \quad \forall i,j\in \mathcal{N};\; t\in[T]\\[4pt]
    &\tilde{m}_{jk}^t 
        = \sum_{i\in \mathcal{N}} x_{ijk}^t,
        &\quad \forall j,k\in \mathcal{N};\; t\in[T]\label{tilde m def}\\[4pt]
    &\tilde{m}_{jk}^{t} \leq M_{jk}^{t}, &\quad \forall j,k\in \mathcal{N}; t\in [T]\label{PostConstraint}\\
    & \sum_{i':d_{i'j}\le d_{ij}} x^t_{i'jk} \geq {\color{red}\delta_{ijk}^t \, M_{jk}^t}, &\forall i,j,k\in \mathcal{N},t\in [T] \label{deltaM}\\
    & x_{ijk}^t \le Q_1 \,\eta_{ijk}^t & \forall i,j,k\in \mathcal{N},t\in [T] \label{Qeta}\\
    & p_{jk} - d_{ij} - c_{jk} \ge s_i^t - Q_2\,(1-\eta^t_{ijk}) & \forall i,j,k\in \mathcal{N},t\in[T] \label{si as lower bound}\\
    & \sum_{j,k\in\mathcal{N}}x^t_{ijk} \ge W_i^t - Q_1 (1-\zeta_i^t) 
    & \forall i\in\mathcal{N},t\in[t]\label{is there lazy worker}\\
    & s_i^t \le  \zeta_i^t Q_2 & \forall i\in\mathcal{N}, t\in[T] \label{stability end}\\
    &s_i^t \ge p_{jk} - d_{ij} - c_{jk} - \delta_{ijk}^t\,Q_2 & \forall i,j,k\in \mathcal{N},t\in [T]\label{si bigger if not full}\\
    & p_{ij},\,
    Y_i^t,\,
    A_i^t,\,
    Y_{ij}^t,\,
    L_i^t,\,
    F_j^t,\,
    \bar{F}_j^t,\,
    U_i^t,\,
    W_i^t,\,
    M_{jk}^t,\,
    \tilde{m}_{jk}^t,\,
    s_i
    \ge 0 & \nonumber\\[4pt]
    &x_{ijk}^t, \hat{m}_{ij}^t\in \mathbb{N},\,
    \delta_{ijk}^t,\,
    \zeta_i^t,\,
    \eta_{ijk}^t,\,
    \in \{0,1\}, \,
    & \quad \forall i,j,k\in \mathcal{N}; t\in [T]
\end{align}
\end{subequations}
\end{itemize}


\section{Reformulation and Solving}

\subsection{Dealing with bilinear terms}

\begin{itemize} 
    \item $ p_{ij}\ \tilde{m}_{ij}^t$ in objective function with continuous variable product with integer variable. Modern commercial solvers such as Gurobi can now efficiently handle non-convex bilinear optimization problems. They achieve global optimality by leveraging advanced techniques, most notably McCormick envelope relaxations integrated within a spatial branch-and-bound framework.\citep{gurobi_manual_current, gurobi_miqcp_2020}
    \item $\alpha A_{i}^t$ in equation \eqref{min(A,D)start} to \eqref{min(A,D)}:
    
    Use Big-M reformulation instead of linear combination \citep{glover1975improved}, which will help us get rid of bilinear term:
\begin{subequations}
    \begin{align}
    & Y_i^t \leq A_{i}^t & \forall i\in \mathcal{N},  t\in [T]\label{2a}\\[4pt]
    &  Y_i^t \leq D_{i}^t & \forall i\in \mathcal{N}; t\in [T]\\[4pt]
    & Y_i^t \geq  A_i^t - Q_3 (1-\beta_i^t) & \forall i\in \mathcal{N},t\in [T]\\
    & Y_i^t \geq  D_i^t - Q_3 \beta_i^t & \forall i\in \mathcal{N},t\in [T]\\
    & \beta_i^t \in \{0,1\} & \forall i\in \mathcal{N}, t\in [T]\label{2e}
    \end{align}
\end{subequations}
    \item $\delta_{ijk}^tM_{jk}^t$ in equation (\ref{deltaM}):
    
   We introduce a new variable $v_{ijk}^t$ to substitute the bilinear term $\delta_{ijk}^tM_{jk}^t$ and use similar method as above, \eqref{deltaM} is equivalent to the following constraints:
\begin{subequations}
    \begin{align}
         v_{ijk}^t&\leq \sum_{i':d_{i'j}\leq d_{ij}}x_{i'jk}^t &\forall i,j,k\in \mathcal{N}; t\in [T] \label{3a}\\
        v_{ijk}^t &\leq M_{jk}^t &\forall i,j,k\in \mathcal{N}; t\in [T] \\
        v_{ijk}^t &\leq Q_3 \cdot \delta_{ijk}^t &\forall i,j,k\in \mathcal{N}; t\in [T] \\
        v_{ijk}^t &\geq M_{jk}^t - Q_3 \cdot (1 - \delta_{ijk}^t) &\forall i,j,k\in \mathcal{N}; t\in [T] \\
        v_{ijk}^t &\geq 0 &\forall i,j,k\in \mathcal{N}; t\in [T]\label{3e}
    \end{align}
\end{subequations}
    \item $Q_1$, $Q_2$ and $Q_3$ are big numbers, $Q_1$ is the number of workers, $Q_2$ is a upper bound of net reward for a worker, and $Q_3$ is the number of bikes. That is
    \begin{align*}
        Q_1 & =|\mathcal{W}^0|\\
        Q_2 & = \max_{jk}\{p_{jk}\}\\
        Q_3 & =\sum_{i\in \mathcal{N}}(A^0_i+U^0_i)
    \end{align*}
\end{itemize}


The Problem is then equivalent to 
\begin{align*}
    \mathbf{\max} \quad
    &\eqref{obj}\\[4pt]
    \text{s.t.} \quad & \eqref{OrderFlow}-\eqref{PostConstraint}, \eqref{Qeta}-\eqref{stability end}\\
    & \eqref{2a}-\eqref{2e}\\
    & \eqref{3a}-\eqref{3e}\\
    & p_{ij},\,
    Y_i^t,\,
    A_i^t,\,
    Y_{ij}^t,\,
    L_i^t,\,
    F_j^t,\,
    \bar{F}_j^t,\,
    U_i^t,\,
    W_i^t,\,
    M_{jk}^t,\,
    \tilde{m}_{jk}^t,\,
    s_i,\,
    v_{ijk}^t
    \ge 0 & \nonumber\\[4pt]
    &x_{ijk}^t, \hat{m}_{ij}^t\in \mathbb{N},\,
    \delta_{ijk}^t,\,
    \zeta_i^t,\,
    \eta_{ijk}^t,\,
    \beta_i^t\,
    \in \{0,1\}, \,
    & \quad \forall i,j,k\in \mathcal{N}; t\in [T]
\end{align*}

\subsection{Equivalence to stability constraints}

Slightly different with the definition of stability in  \citep{baiou2000stable}, we define matching and blocking pairs and stability in our problem as follows:

\paragraph{Definition of Matching}: At time $t$, given a locations set $\mathcal{N}^t$ (with the workers $\mathcal{W}^t$ located on these nodes, and $W_i^t>0$ workers at all nodes $i$) and a jobs set $\mathcal{M}^t$, a matching is defined as an array $X^t = \{x_{ijk}^t\}_{i\in \mathcal{N}^t,jk\in \mathcal{M}^t}$, where $x_{ijk}^t$ denotes the number of workers from location $i$ assigned to job $jk$. The induced matching correspondence $\mu^t$ is defined as
\begin{align*}
    & \mu^t(i) = \{jk \in \mathcal{M}^t \mid x_{ijk}^t > 0\} \quad \forall i\in \mathcal{N}^t, \\
    & \mu^t(jk) = \{i \in \mathcal{N}^t \mid x^t_{ijk} > 0\} \quad \forall jk\in \mathcal{M}^t.
\end{align*}

Given a matching defined by $\mu^t$, we define $\hat{s}^t_i$ to be the lowest reward that a worker at $i$ get in the matching, that is,
\begin{align}
    \hat{s}_i^t =
    \begin{cases}
    \min_{jk\in \mu^t(i)} \left(p_{jk}-d_{ij}-c_{jk}\right) & \sum_{jk}x_{ijk}^t = W_i^t,\\
    \min \left\{\min_{jk\in \mu^t(i)} \left(p_{jk}-d_{ij}-c_{jk}\right),0\right\} & \text{o.w.},
    \end{cases}
    \quad \forall i\in \mathcal{N}^t \label{definition s hat}.
\end{align}

\paragraph{Definition of Blocking Pair:} A pair of a worker location $i$ and a job $jk$ at time $t$ is a \textit{blocking pair} if both of the following conditions hold simultaneously:
\begin{enumerate}
    \item \textbf{Worker-side incentive to deviate:} There exists at least one worker at location $i$ who strictly prefers job $jk$ over their current assignment (or their unemployment state). 
    
    Because workers at location $i$ are homogeneous, this implies that the minimum realized utility among all workers at $i$, $\hat{s}_i^t$, is strictly less than the utility provided by job $jk$:
    \begin{equation*}
        \hat{s}_i^t < p_{jk} - d_{ij} - c_{jk}.
    \end{equation*}

    \item \textbf{Job-side incentive to deviate:} Job $jk$ either has unfilled capacity, or it has accepted workers from locations strictly worse (i.e., farther) than $i$. Mathematically, this requires:
    \begin{equation*}
        \sum_{i': d_{i'j} \le d_{ij}} x_{i'jk}^t < M_{jk}^t.
    \end{equation*}
\end{enumerate}

\paragraph{Definition of Stability:} A matching given by the aggregate flow $x_{ijk}^t$ is \textbf{stable} if and only if no such blocking pair $(i,jk)$ exists in the network, i.e., the two "blocking pair conditions" do not hold simultaneously. 

Now we show that mathematically, the definition above is equivalent to satisfying the constraints in the model.

\begin{lemma}[Equivalence of Stability and Reformulation]
    A matching given by the aggregate flow $x_{ijk}^t$ is stable if and only if there exist auxiliary variables $(\delta_{ijk}^t, \eta_{ijk}^t, \zeta_i^t, s_i^t)$ such that the set of constraints in the model is satisfied.
\end{lemma}

\paragraph{Necessity}: We first show that given a stable matching $X=\left\{x_{ijk}^t\right\}_{i\in\mathcal{N}^t,jk\in \mathcal{M}^t}$, whose induced correspondence is $\mu$, we can construct auxiliary variables $\delta_{ijk}^t$, $\eta_{ijk}^t$, $\zeta_i^t$ and $s_i^t$, such that the decision variables $\left\{x_{ijk}^t,\delta_{ijk}^t, \eta_{ijk}^t, \zeta_i^t, s_i^t\right\}$ satisfy the constraints.

We construct the auxiliary variables for all $i,j,k$ as follows:
\begin{align*}
    & \delta_{ijk}^t = \begin{cases}
        1 & \sum_{i':d_{i'j}\le d_{ij}} x^t_{i'jk} = M_{jk}^t\\
        0 & \text{o.w.}
    \end{cases}\\
    & \eta_{ijk}^t = \begin{cases}
        1 & x_{ijk}^t > 0\\
        0 & \text{o.w.}
    \end{cases}\\
    & \zeta_i^t = \begin{cases}
        1 & \sum_{jk\in \mathcal{M}^t}x_{ijk}^t=W_i^t\\
        0 & \text{o.w.}
    \end{cases}\\
    & s_i^t = \hat{s}_i^t,
\end{align*}
and check whether the constraints are satisfied.
\begin{itemize}
    \item \eqref{deltaM}: By our construction of $\delta_{ijk}^t$, either of the following two cases holds: 
    if $\sum_{i':d_{i'j}\le d_{ij}} x^t_{i'jk} = M_{jk}^t$ and we construct $\delta_{ijk}^t=1$, thus \eqref{deltaM} becomes $\sum_{i':d_{i'j}\le d_{ij}} x^t_{i'jk} \geq M_{jk}^t$ and holds; 
    if $\sum_{i':d_{i'j}\le d_{ij}} x^t_{i'jk} < M_{jk}^t$ and we construct $\delta_{ijk}^t=0$, thus \eqref{deltaM} becomes $\sum_{i':d_{i'j}\le d_{ij}} x^t_{i'jk} \geq 0$ and holds;
    \item \eqref{Qeta}: By our construction of $\eta_{ijk}^t$, either of the following two cases holds: 
    if $x_{ijk}^t=0$ and we construct $\eta_{ijk}^t=0$, thus \eqref{Qeta} becomes $x_{ijk}^t\leq 0$ and holds; 
    if $x_{ijk}^t>0$ and we construct $\eta_{ijk}^t=1$, thus \eqref{Qeta} becomes $x_{ijk}^t \leq Q_1$ ($Q_1$ is a big number) and holds;
    \item \eqref{si as lower bound}: By our construction
    \begin{align*}
        s_i^t = \begin{cases}
    \min_{jk\in \mu(i)} \left(p_{jk}-d_{ij}-c_{jk}\right) & \sum_{jk}x_{ijk} = W_i,\\
    \min \left\{\min_{jk\in \mu(i)} \left(p_{jk}-d_{ij}-c_{jk}\right),0\right\} & \text{o.w.},
    \end{cases}
    \quad \forall i\in \mathcal{N}.
    \end{align*} Either of the following two cases holds: 
    if $x_{ijk}^t >0$, we construct $\eta_{ijk}^t=1$ which implies that $jk\in \mu(i)$, \eqref{si as lower bound} becomes $s_i^t\leq p_{jk}-d_{ij}-c_{jk}$ and holds by our construction of $s_i^t$; 
    if $x_{ijk}^t =0$, we construct $\eta_{ijk}^t=0$ which implies that $jk\notin \mu(i)$, \eqref{si as lower bound} becomes $s_i^t\leq p_{jk}-d_{ij}-c_{jk} + Q_2$ and holds for big number $Q_2$ by our construction of $s_i^t$;
    \item \eqref{is there lazy worker}: By our construction of $\zeta_{i}^t$, either of the following two cases holds: 
    if $\sum_{jk\in\mathcal{M}^t} x_{ijk}^t=W_i^t$ and we construct $\zeta_{i}^t=1$, thus \eqref{is there lazy worker} becomes $\sum_{jk\in\mathcal{M}^t}x_{ijk}^t\geq W_i^t$ and holds;
    if $\sum_{jk\in\mathcal{M}^t} x_{ijk}^t<W_i^t$ and we construct $\zeta_{i}^t=0$, thus \eqref{is there lazy worker} becomes $\sum_{jk\in\mathcal{M}^t}x_{ijk}^t\geq W_i^t - Q_1$ and holds for big number $Q_1$;
    \item \eqref{stability end}: By our construction, either of the following two cases holds: 
    if $\sum_{jk\in\mathcal{M}^t} x_{ijk}^t=W_i^t$ and we construct $\zeta_{i}^t=1$, thus \eqref{stability end} becomes $s_i^t\leq Q_2$ and holds for big number $Q_2$;
    if $\sum_{jk\in\mathcal{M}^t} x_{ijk}^t<W_i^t$, we construct $\zeta_{i}^t=0$ and $s_i^t=\min \left\{\min_{jk\in \mu(i)} \left(p_{jk}-d_{ij}-c_{jk}\right),0\right\}$, \eqref{stability end} becomes $s_i^t\leq 0$ and holds by our construction of $s_i^t$;
    \item \eqref{si bigger if not full}: Either of the following cases holds: 
    if $\sum_{i':d_{i'j}\le d_{ij}} x^t_{i'jk} = M_{jk}^t$, we construct $\delta_{ijk}^t=1$, \eqref{si bigger if not full} becomes $s_i^t \leq p_{jk} - d_{ij} - c_{jk} - Q_2$ and holds for big number $Q_2$; 
    if $\sum_{i':d_{i'j}\le d_{ij}} x^t_{i'jk} < M_{jk}^t$, we construct $\delta_{ijk}^t=0$, \eqref{si bigger if not full} becomes $s_i^t \geq p_{jk} - d_{ij} - c_{jk}$ and holds, because the matching is stable, which implies that $\sum_{i':d_{i'j}\le d_{ij}} x^t_{i'jk} < M_{jk}^t$ and $s_i^t=\hat{s}_i^t<p_{ij}-d_{jk}-c_{jk}$ does not hold simultaneously, thus we must have $s_i^t \geq p_{jk} - d_{ij} - c_{jk}$.
\end{itemize}


\paragraph{Sufficiency}: We then show that a job assignment flow $\left\{x_{ijk}^t\right\}_{i,j,k}$ that satisfies constraints from the model is a stable matching.

To show that the matching is stable, we only need to show that for all $(i,jk)$ pair, the two conditions do not hold simultaneously. That is,
\[
\sum_{i': d_{i'j} \le d_{ij}} x_{i'jk}^t < M_{jk}^t
\]
and 
\[
\hat{s}_i^t < p_{jk} - d_{ij} - c_{jk}
\]
does not hold simultaneously for all $i,j,k$. We consider two cases based on the value of $\delta^t_{ijk}$ for any $i,j,k$.

For those $i,j,k$ such that $\sum_{i': d_{i'j} \le d_{ij}} x_{i'jk}^t = M_{jk}^t$, the first condition is violated and the two conditions must not hold simultaneously, i.e., the pair is not a blocking pair.

For those $i,j,k$ such that $\sum_{i': d_{i'j} \le d_{ij}} x_{i'jk}^t < M_{jk}^t$, we must show that 
\begin{align}\label{eq:sihatt}
\hat{s}_i^t \geq p_{jk} - d_{ij} - c_{jk}.
\end{align}

Since \eqref{deltaM} holds and $\sum_{i': d_{i'j} \le d_{ij}} x_{i'jk}^t < M_{jk}^t$, $\delta_{ijk}^t$ must be 0, and \eqref{si bigger if not full} becomes
\[
s_i^t \ge p_{jk} - d_{ij} - c_{jk}. 
\]
Therefore, to prove \eqref{eq:sihatt}, it is sufficient to prove that $s_i^t \leq \hat{s}_i^t$. We consider two cases separately:
\begin{itemize}
    \item If $\sum_{jk\in\mathcal{M}^t}x_{ijk}^t=W_i^t$, by the definition of $\hat{s}_i^t$ in \eqref{definition s hat}, we have
    \[
    \hat{s}_i^t = \min_{jk\in \mu^t(i)} \left(p_{jk}-d_{ij}-c_{jk}\right).
    \]
    For any job $j'k'\in \mu^t(i)$, $x_{ij'k'}^t>0$ and since \eqref{Qeta} holds, $\eta_{ij'k'}^t$ must be $1$ and \eqref{si as lower bound} becomes 
    \[
    s_i^t \leq p_{j'k'} - c_{j'k'} - d_{ij'} \quad \forall j'k' \in \mu^t(i).
    \]
    Thus we have 
    \[
    s_i^t \leq \min_{jk\in \mu^t(i)} \left(p_{jk}-d_{ij}-c_{jk}\right) = \hat{s}_i^t.
    \]

    \item If $\sum_{jk\in\mathcal{M}^t}x_{ijk}^t<W_i^t$, since \eqref{is there lazy worker} holds, $\zeta_i^t$ must be $0$; and by the definition of $\hat{s}_i^t$ in \eqref{definition s hat}, we have
    \[
    \hat{s}_i^t = \min\left\{0, \min_{jk\in \mu^t(i)} \left(p_{jk}-d_{ij}-c_{jk}\right)\right\}.
    \]
    For those jobs $j'k'\in \mu^t(i)$, $x_{ij'k'}^t>0$ and since \eqref{Qeta} holds, $\eta_{ij'k'}^t$ must be $1$ and \eqref{si as lower bound} becomes 
    \[
    s_i^t \leq p_{jk} - c_{jk} - d_{ij} \quad \forall i,j,k,
    \]
    Thus we have 
    \[
    s_i^t \leq \min_{jk\in \mu^t(i)} \left(p_{jk}-d_{ij}-c_{jk}\right).
    \]
    Additionally, since $\zeta_i^t=0$, from \eqref{stability end} we know that $s_i^t \leq 0$. Thus, combining the two results, we know that
    \[
    s_i^t \leq \min\left\{0, \min_{jk\in \mu^t(i)} \left(p_{jk}-d_{ij}-c_{jk}\right)\right\} = \hat{s}_i^t.
    \]
\end{itemize}


\newpage

\subsection{Stochastic Dual Dynamic Integer Programming\citep{zou2019stochastic}}\label{sec:sddip}


Formally, SDDiP targets problems of the form
\begin{equation}\label{eq:msip-standard}
    Q_n(x_{a(n)}) = \min_{x_n,\,y_n}\left\{
    f_n(x_n,y_n) + \sum_{m\in C(n)} q_{nm}\,Q_m(x_n)
    \;:\;
    (x_{a(n)},x_n,y_n)\in X_n,\;
    x_n\in\{0,1\}^d
    \right\},
\end{equation}
where $x_n\in\{0,1\}^d$ is the binary state variable linking node $n$ to its children, $y_n$ is a local (possibly mixed-integer) variable, $f_n$ is a linear objective, and $X_n$ is a compact mixed-integer polyhedral set. For problems with general integer or continuous state variables, \citep{zou2019stochastic} show that they can be approximated to precision $\varepsilon$ via binary expansion, with the number of binary state variables per stage bounded by $k \le d(\lfloor\log_2(M/\varepsilon)\rfloor+1)$, where $d$ is the original state dimension and $M$ depends on problem data.

We now examine the structural requirements for applying SDDiP to our micromobility operations problem, and identify the challenges that arise.


\subsubsection{Requirement 1: MILP Stage Subproblems}\label{sec:milp-req}

SDDiP requires that each stage subproblem is a mixed-integer linear program (MILP).Bilinear term $p_{jk}\,\tilde{m}_{jk}^t$ in the objective.



\subsubsection{Requirement 2: Binary State Variables and DP Reformulation}\label{sec:binary-state}

SDDiP requires the state variables---those linking consecutive stages---to be binary. Applying this to our problem involves two steps: (i)~restoring the Markov property via pipeline augmentation, and (ii)~binarizing all integer state variables.

\paragraph{Step 1: Pipeline augmentation.}
The state transitions \eqref{FlowTransition} involve multi-period delays: user trips from $i$ to $j$ require $t_{ij}$ periods, and worker tasks from $i$ via $(j,k)$ require $\delta_{ijk}:=d_{ij}+c_{jk}$ periods. The delayed terms $Y_{ij}^{t-t_{ij}}$ and $x_{ijk}^{t-\delta_{ijk}}$ make the system non-Markovian. We restore the Markov property by introducing \emph{pipeline state variables} that track all in-transit entities.

For each OD pair $(i,j)$ with travel time $t_{ij}\ge 2$, define
\[
    P_{ij,r}^t := \text{number of bikes in transit from $i$ to $j$, arriving in $r$ more periods},
    \quad r=1,\ldots,t_{ij}-1.
\]
These evolve via
\begin{align}
    P_{ij,r}^{t+1} &= P_{ij,r+1}^{t},
    &r=1,\ldots,t_{ij}-2, \label{eq:pipe-bike}\\
    P_{ij,t_{ij}-1}^{t+1} &= Y_{ij}^{t+1}. \nonumber
\end{align}
The arrival terms then become functions of the current state: $F_j^t = \sum_i P_{ij,1}^t(1-\phi_{ij})$ and $\bar{F}_j^t = \sum_i P_{ij,1}^t\phi_{ij}$.

Similarly, for each worker-task triple $(i,j,k)$ with delay $\delta_{ijk}\ge 2$, define
\[
    G_{ijk,r}^t := \text{number of workers dispatched from $i$ for task $(j,k)$, arriving at $k$ in $r$ more periods},
\]
for $r=1,\ldots,\delta_{ijk}-1$, with an analogous shift-and-entry transition:
\begin{align}
    G_{ijk,r}^{t+1} &= G_{ijk,r+1}^{t},
    &r=1,\ldots,\delta_{ijk}-2, \label{eq:pipe-worker}\\
    G_{ijk,\delta_{ijk}-1}^{t+1} &= x_{ijk}^{t+1}. \nonumber
\end{align}

The complete Markov state at time $t$ (excluding pricing) is then
\begin{equation}\label{eq:state-op}
\boxed{
    \mathbf{s}_t^{\mathrm{op}} = \Big(\;
    \underbrace{A_j^t,\; U_j^t,\; W_j^t}_{j\in\mathcal{N}},\;\;
    \underbrace{M_{jk}^t}_{j,k\in\mathcal{N}},\;\;
    \underbrace{P_{ij,r}^t}_{\substack{i,j\in\mathcal{N}\\r=1,\ldots,t_{ij}-1}},\;\;
    \underbrace{G_{ijk,r}^t}_{\substack{i,j,k\in\mathcal{N}\\r=1,\ldots,\delta_{ijk}-1}}
    \;\Big).
}
\end{equation}

Using these pipeline variables, the original transition equations \eqref{FlowTransition} are rewritten as
\begin{subequations}\label{eq:transitions-rewrite}
\begin{align}
    A_j^{t+1}
    &= A_j^t - Y_j^t
    + \textstyle\sum_{i} P_{ij,1}^t(1-\phi_{ij})
    - \sum_{k\neq j} \hat{m}_{jk}^t
    + \sum_{i,k} G_{ikj,1}^t,
    &\quad \forall j, \label{eq:A-trans}\\[3pt]
    U_j^{t+1}
    &= U_j^t
    + \textstyle\sum_{i} P_{ij,1}^t\phi_{ij}
    - \sum_{i} G_{ijj,1}^t,
    &\quad \forall j, \\[3pt]
    W_k^{t+1}
    &= W_k^t
    - \textstyle\sum_{i,j} x_{kij}^t
    + \sum_{i,j} G_{ijk,1}^t,
    &\quad \forall k, \\[3pt]
    M_{jk}^{t+1}
    &= M_{jk}^t - \tilde{m}_{jk}^t + \hat{m}_{jk}^{t+1},
    &\quad \forall j,k. \label{eq:M-trans}
\end{align}
\end{subequations}
All arrival terms now depend only on the current state $\mathbf{s}_t^{\mathrm{op}}$, restoring the Markov property.

\paragraph{Step 2: Binarization.}
All components of $\mathbf{s}_t^{\mathrm{op}}$ are nonneg.\ integers with known upper bounds. By Theorem~3 of \citep{zou2019stochastic}, each integer variable $s\in\{0,1,\ldots,U\}$ is replaced by its binary expansion $s=\sum_{l=1}^\kappa 2^{l-1}\lambda_l$ with $\lambda_l\in\{0,1\}$ and $\kappa=\lfloor\log_2 U\rfloor+1$. Table~\ref{tab:binary-op} counts the resulting binary state variables for a representative instance.

\begin{table}[htbp]
\centering
\caption{Binary state variable count per stage (operational state only) for $n=|\mathcal{N}|=10$, $B_{\max}=100$, $W=30$, $M_{\max}=20$, uniform $t_{ij}=3$, uniform $\delta_{ijk}=5$.}
\label{tab:binary-op}
\small
\begin{tabular}{l r r r r}
\hline
\textbf{Variable group} & \textbf{Integer vars} & $\kappa$ &
  \textbf{Binary vars} & \textbf{\% of total}\\
\hline
$A_j^t$       & 10    & 7 & 70       & 0.3\% \\
$U_j^t$       & 10    & 7 & 70       & 0.3\% \\
$W_k^t$       & 10    & 5 & 50       & 0.2\% \\
$M_{jk}^t$    & 100   & 5 & 500      & 2.3\% \\
$P_{ij,r}^t$  & $100\times 2=200$  & 7 & 1{,}400   & 6.3\% \\
$G_{ijk,r}^t$ & $1{,}000\times 4=4{,}000$ & 5 & 20{,}000  & 90.6\% \\
\hline
\textbf{Total $K_{\mathrm{op}}$}
              &       &   & \textbf{22{,}090} & 100\% \\
\hline
\end{tabular}
\end{table}

The worker pipeline $G_{ijk,r}^t$ dominates at $90.6\%$ of the total, scaling as $O(n^3\bar{\delta})$ where $\bar{\delta}$ is the average worker travel delay.

\paragraph{Comparison with existing SDDiP benchmarks.}
As a reference point, the portfolio optimization experiment in \citep{zou2019stochastic} involves $n=100$ stocks from the S\&P~100 index with continuous state variables (asset values) approximated via binary expansion to precision $\varepsilon=10^{-2}$, resulting in approximately 1{,}500 binary state variables per stage. That experiment achieves optimality gaps of $1$--$2\%$ within 10 iterations across horizons of $T=4$ to $T=12$ stages. Our operational state of $K_{\mathrm{op}}=22{,}090$ binary variables per stage is roughly 15 times larger, suggesting that directly applying SDDiP to the full model (with all pipeline variables) would face significant computational challenges. Dimensionality reduction---via instant-travel approximations ($t_{ij}=\delta_{ijk}=1$, eliminating all pipeline variables and reducing $K_{\mathrm{op}}$ to 690), network aggregation ($n=5$ yields $K_{\mathrm{op}}\approx 3{,}300$), or sparsification of inactive OD triples---would be essential for computational tractability.




\subsection{Progressive Hedging Algorithm\citep{gade2016obtaining}}

Denote the pricing decision as $p$, the job posting decision as $m$ and all other decision variables as $x$, since the constraints are all linear after reformulation, the problem can be modeled as
\begin{align}
    \mathbf{\max} \quad
    &\mathbb{E}\left[
    \sum_{t=1}^{T}
    \left(
    \sum_{i\in \mathcal{N}} \big( \sum_{j\in \mathcal{N}}R_{ij}\, Y_{ij}^t - C_p\, L_{i}^t \big)
    - \sum_{j,k\in \mathcal{N}} p_{jk}\, \tilde{m}_{jk}^t
    \right)
    \right]\\[4pt]
    \textbf{s.t.}\nonumber\\
    & A_p\,p + A_m\, m + A_x\,x \le b.
\end{align}
Suppose the stochasticity is modeled as several scenarios, each scenario is denoted as $s\in \mathcal{S}$, then the model can be written as (denote the objective function as $f(p,m,x)$):

\begin{align}
    \mathbf{\max} \quad
    &\sum_{s\in \mathcal{S}}
    f\big(p(s),m(s),x(s)\big)\pi_s\\[4pt]
    \textbf{s.t.}\nonumber\\
    & A_p(s)\,p(s) + A_m(s)\, m(s) + A_x(s)\,x(s) \le b(s) \label{constraints} & \forall s\in \mathcal{S}\\
    &p(s)=p(s')=\hat{p} & \forall s,s'\in \mathcal{S}.\label{non-anticipativity constraint}
\end{align}

Solve this problem by solving
\begin{align}
    \mathbf{\max} \quad
    &\sum_{s\in \mathcal{S}} \left[
    f\big(p(s),m(s),x(s)\big)\pi_s -\lambda_s(p(s)-\hat{p}) + {\rho\over2}(p(s)-\hat{p})^2
    \right]\\[4pt]
    \textbf{s.t.}\nonumber\\
    & A_p(s)\,p(s) + A_m(s)\, m(s) + A_x(s)\,x(s) \le b \label{constraints} & \forall s\in \mathcal{S}.
\end{align}

The objective can be decomposed into several scenarios, and we can solve this by using the following procedure
\begin{itemize}
    \item For iteration $k=1,2,...,N$:
    \begin{itemize}
        \item fix $\hat{p}^{k}$;
        \item For scenarios $s=1,2,...,S$:
        \begin{itemize}
            \item Solve deterministic subproblem
            \begin{align*}z^*(s)&=f(p^*(s),m^*(s),x^*(s))\\
            &=\max_{p(s),m(s),x(s)}
            \left\{
            f(p(s),m(s),x(s))\mid A_p(s)\,p(s) + A_m(s)\, m(s) + A_x(s)\,x(s) \le b(s)
            \right\}
            \end{align*}
        \end{itemize}
        \item Update $\hat{p}^{k+1}=\sum_{s\in \mathcal{S}}\pi_s p^*(s)$
        \item Update $\lambda_s^{k+1} = \lambda_s^{k} + \rho\pi_s (p^*(s)-\hat{p}^{k+1})$
    \end{itemize}
\end{itemize}

\paragraph{Multistage Problems} For a multi-stage stochastic problem, let $\xi^{[t]}(s) = 
(\xi_1(s),\ldots,\xi_t(s))$ denote the partial history of scenario 
$s$ up to stage $t$. Two scenarios $s$ and $s'$ are 
\emph{indistinguishable} at stage $t$ if $\xi^{[t]}(s)=\xi^{[t]}(s')$; 
this defines an equivalence relation whose classes are the 
\emph{scenario tree nodes} at stage $t$. Let $\mathcal{T}_t$ denote 
the set of all nodes at stage $t$, with each node $n\in\mathcal{T}_t$ 
identified by its partial history $\xi^{[t]}$.

We then define:
\begin{itemize}
    \item $n(s,t)\in\mathcal{T}_t$: the node at stage $t$ to which 
          scenario $s$ belongs, i.e.\ $n(s,t)=\xi^{[t]}(s)$;
    \item $\mathcal{S}(n)=\{s\in\mathcal{S}: n(s,t)=n\}$: the set 
          of all scenarios passing through node $n\in\mathcal{T}_t$;
    \item $\pi_n = \sum_{s\in\mathcal{S}(n)}\pi_s$: the marginal 
          probability of reaching node $n$.
\end{itemize}
The non-anticipativity constraints (NAC) then require that decisions 
depend only on the information available at each node:
\begin{align}
    p(s) &= \hat{p} 
    & \forall s\in\mathcal{S}, \label{nac-p}\\
    m^t(s) &= \hat{m}^t_{n(s,t)} 
    & \forall s\in\mathcal{S},\; t=1,\ldots,T. \label{nac-m}
\end{align}
Relaxing \eqref{nac-p}--\eqref{nac-m} via an augmented Lagrangian 
yields the decomposed subproblem for scenario $s$ at iteration $k$:
\begin{align}
    z^*(s) = \max_{p(s),m(s),x(s)} \Bigg\{
    &\pi_s f\big(p(s),m(s),x(s)\big) \notag\\
    &-\lambda^{0,k}_s\big(p(s)-\hat{p}^k\big)
     +\tfrac{\rho}{2}\big(p(s)-\hat{p}^k\big)^2 \notag\\
    &-\sum_{t=1}^{T}\!\left[
        \lambda^{t,k}_s\big(m^t(s)-\hat{m}^{t,k}_{n(s,t)}\big)
       +\tfrac{\rho}{2}\big(m^t(s)-\hat{m}^{t,k}_{n(s,t)}\big)^2
      \right] \Bigg\} \notag\\
    \text{s.t.}\quad
    &A_p(s)\,p(s)+A_m(s)\,m(s)+A_x(s)\,x(s)\le b(s).
\end{align}
The algorithm proceeds as follows. For iteration $k=1,2,\ldots$:
\begin{itemize}
    \item Fix $\hat{p}^k$ and $\hat{m}^{t,k}_n$ for all $t,n$.
    \item For each $s\in\mathcal{S}$ (in parallel): solve the 
          deterministic subproblem to obtain 
          $p^*(s),\{m^{t*}(s)\},x^*(s)$.
    \item Update consensus:
    \begin{align}
        \hat{p}^{k+1} &= \sum_{s\in\mathcal{S}}\pi_s\,p^*(s),\\
        \hat{m}^{t,k+1}_n &= \frac{1}{\pi_n}
            \sum_{s\in\mathcal{S}(n)}\pi_s\,m^{t*}(s)
            \quad \forall\,t,\,\forall\,n\in\mathcal{T}_t.
    \end{align}
    \item Update multipliers for all $s\in\mathcal{S}$:
    \begin{align}
        \lambda^{0,k+1}_s &= \lambda^{0,k}_s
            +\rho\,\pi_s\big(p^*(s)-\hat{p}^{k+1}\big),\\
        \lambda^{t,k+1}_s &= \lambda^{t,k}_s
            +\rho\,\pi_s\big(m^{t*}(s)
            -\hat{m}^{t,k+1}_{n(s,t)}\big)
            \quad \forall\,t.
    \end{align}
    \item Terminate if $|p^*(s)-\hat{p}^{k+1}|\le\varepsilon$ and 
          $|m^{t*}(s)-\hat{m}^{t,k+1}_{n(s,t)}|\le\varepsilon$ 
          for all $s,t$.
\end{itemize}


\printbibliography
\end{document}